%% ============================================================
%%  Chapitre 4 — Architecture de la Plateforme
%% ============================================================
\chapter{Architecture de la Plateforme}
\label{ch:architecture}

\section{Paradigme Méta-Stratégie~: les Quatre Stations}
\label{sec:quatre_stations}

\subsection{Séparation des responsabilités}

L'architecture de la plateforme découle directement du paradigme méta-stratégie de
López de Prado~\cite{lopezdeprado2018}. Quatre stations spécialisées et indépendantes
constituent la ligne de production quantitative~:

\begin{quote}
\textit{«~Backtesting is not a research tool. Feature importance is.~»}
— López de Prado (2018), Ch.~7.
\end{quote}

Cette séparation n'est pas cosmétique~: elle prévient l'overfitting cognitif qui
survient quand le chercheur observe les résultats du backtest et \emph{modifie} en
conséquence les paramètres du signal — une forme de fuite d'information entre stations.

\begin{table}[H]
\centering
\caption{Les quatre stations de la plateforme}
\label{tab:quatre_stations}
\begin{tabularx}{\linewidth}{@{}lllY@{}}
\toprule
\textbf{Station} & \textbf{Page}    & \textbf{Route}     & \textbf{Mandat} \\
\midrule
Data Curators    & Données          & \texttt{/data}     & Gérer l'univers de marché OHLCV : ingestion, stockage, fraîcheur \\
Feature Analysts & Signaux          & \texttt{/signals}  & Découvrir quels indicateurs ont un pouvoir prédictif (OOS) \\
Strategists      & Stratégie        & \texttt{/strategy} & Définir les règles de trading et paramétrer pour WFO \\
Backtesting      & Backtest         & \texttt{/backtest} & Évaluer les stratégies via WFA, produire un verdict de robustesse \\
\bottomrule
\end{tabularx}
\end{table}

\section{Stack Technologique}
\label{sec:stack}

La migration d'un prototype Streamlit vers une architecture full-stack de production
a été réalisée en Semaine~5~\cite{rapport_s5, rapport_point_hebdo}. Le choix de la
stack repose sur quatre critères~: scalabilité des calculs, expérience utilisateur
fluide, traçabilité des runs et reproductibilité.

\begin{figure}[H]
\centering
\placeholder{Diagramme d'architecture complet — style "stacké de gauche à droite" :\\
\textbf{Frontend (Next.js 14)} : App Router, Tailwind CSS, Recharts, SWR polling 30s. \\
\textbf{API (FastAPI)} : Pydantic v2, SQLAlchemy, 18 endpoints market\_data, routers runs/signals/strategy. \\
\textbf{Workers (RQ / Redis)} : Calculs asynchrones, exécution des runs, optimisation parallèle. \\
\textbf{Moteur quant (quant\_core)} : Data → Signal Engine A→G → WFO → BacktestEngine → Artifacts. \\
\textbf{PostgreSQL} : Tables stock\_master, dataset, runs, fills, metrics, strategy\_decision. \\
\textbf{MinIO / S3} : Stockage objet — parquets OHLCV, artefacts plots, exports. \\
Flèches HTTP, SQL, RQ enqueue, S3 GET/PUT. \\
Source : Point hebdomadaire BMCE Capital, Semaine 5.}
\caption{Architecture technique de la plateforme~: stack modulaire et industrialisable.
Chaque composant est isolé et communique via des interfaces définies.
(Adapté de~\cite{rapport_point_hebdo})}
\label{fig:architecture_stack}
\end{figure}

\subsection{Détail des composants}

\begin{table}[H]
\centering
\caption{Composants de la stack et leurs rôles}
\label{tab:stack_detail}
\begin{tabularx}{\linewidth}{@{}lYY@{}}
\toprule
\textbf{Composant}        & \textbf{Technologie}             & \textbf{Rôle clé} \\
\midrule
Frontend                  & Next.js 14 (App Router), Tailwind, SWR & Pages interactives, polling automatique, graphiques \\
API                       & FastAPI, Pydantic v2, SQLAlchemy & Orchestration, validation, 18+ endpoints REST \\
Queue asynchrone          & RQ (Redis Queue)                 & Découpler les runs lourds du thread HTTP \\
Noyau quantitatif         & Python, Numba, Pandas, NumPy     & Signal Engine, WFO, métriques, artefacts \\
Base de données           & PostgreSQL                       & Persistance des runs, métriques, fills, datasets \\
Stockage objet            & MinIO (compatible S3)            & Parquets OHLCV, plots PNG, exports JSON \\
Migrations                & Alembic                          & Versionnage du schéma BDD \\
Tests                     & pytest                           & 87+ tests unitaires et d'intégration \\
\bottomrule
\end{tabularx}
\end{table}

\section{Station 1 --- Données}
\label{sec:station_data}

\subsection{Pipeline d'ingestion OHLCV}

La page Données est le centre de contrôle de l'ensemble de la data pipeline.
Tout signal et tout backtest dépend des données gérées ici.

Les données OHLCV (Open, High, Low, Close, Volume) sont ingérées depuis trois
sources adaptatives~:

\begin{itemize}
    \item \textbf{YFinance Morocco Adapter}~: données gratuites pour les tickers MASI~;
    \item \textbf{Bourse Directe Adapter}~: données officielles de la Bourse de Casablanca~;
    \item \textbf{BDCSession Adapter}~: session authentifiée pour les données institutionnelles.
\end{itemize}

Trois formats de fichiers sont détectés automatiquement (BMCE new/old, Investing.com),
avec gestion robuste des encodages (déaccent, casefold, mojibake) et des nombres
au format français (séparateur décimal virgule).

\subsection{Calendrier des jours fériés}

Un calendrier des jours fériés spécifique à la Bourse de Casablanca est maintenu
pour les 93 tickers MASI. Il distingue deux niveaux de certitude~:
\begin{itemize}
    \item \textbf{Confirmés}~: dates officiellement publiées~;
    \item \textbf{Provisoires}~: estimations basées sur les calendriers religieux et civils.
\end{itemize}

\subsection{Fraîcheur et rafraîchissement automatique}

Un scheduler APScheduler déclenche un rafraîchissement automatique quotidien à
\textbf{18h00, heure d'Afrique/Casablanca}, du lundi au vendredi (après clôture).
Ce comportement est contrôlable via la variable d'environnement
\texttt{MARKET\_REFRESH\_CRON\_ENABLED}.

\subsection{Stockage canonique}

Chaque dataset est stocké sous une clé canonique unique~:

\begin{equation}
\texttt{key} = \texttt{datasets/\{data\_hash\}/\{filename\}}
\label{eq:canonical_key}
\end{equation}

La valeur de \texttt{data\_hash} est calculée au moment de l'upload. Cette clé
est persistée en base de données comme \og source of truth \fg{} — elle n'est
jamais recalculée sauf pour les lignes héritées (migrations Alembic).

\figuretodo{screenshot_data_page}{Capture d'écran de la page Données~: catalogue des 93 tickers MASI avec statut de fraîcheur, date de dernière mise à jour, bouton de rafraîchissement manuel et vue graphique prix.}

\section{Station 2 --- Signaux (Interface Utilisateur)}
\label{sec:station_signaux}

La page Signaux expose le Signal Engine (Chapitre~\ref{ch:signaux}) via une interface
visuelle intuitive.

\subsection{Vue par ticker}

Pour chaque ticker sélectionné, l'utilisateur voit~:

\begin{itemize}
    \item Un \textbf{compteur de vitesse principal} (gauge $[-100, +100]$) représentant
    le score d'ensemble toutes familles confondues~;
    \item \textbf{Quatre sous-compteurs} (Tendance, Momentum, Oscillation, Volume)
    affichant le score par catégorie~;
    \item La \textbf{recommandation directionnelle}~: FORT ACHAT, ACHAT, NEUTRE,
    VENTE, FORTE VENTE.
\end{itemize}

\figuretodo{screenshot_signals_page}{Capture d'écran de la page Signaux~: ticker IAM, compteur de vitesse central (+67, signal ACHAT), 4 sous-compteurs par famille, tableau des variantes représentantes avec labels HAUSSIER / SURVENDU / ACCUMULATION.}

\subsection{Drill-down par variante}

L'utilisateur peut cliquer sur un représentant pour accéder au détail~:
\begin{itemize}
    \item Paramètres de la variante~;
    \item Résultats OOS par fenêtre temporelle (Sharpe, MDD, Win\%)~;
    \item Décomposition du reliability\_score (4 composantes)~;
    \item Graphique du signal dans le temps superposé au prix.
\end{itemize}

\figuretodo{screenshot_variant_detail}{Capture d'écran du panneau de détail d'une variante SMA(50) sur IAM~: OOS par fenêtre (tableau + barplot Sharpe), décomposition reliability score (radar chart), signal timeline.}

\section{Station 3 --- Stratégie}
\label{sec:station_strategie}

La page Stratégie permet de définir une stratégie de trading complète avant
tout backtesting.

\subsection{Structure en cinq blocs}

Conformément à l'architecture définie en Semaine~9~\cite{rapport_s9}~:

\begin{enumerate}
    \item \textbf{Univers}~: sélection des actions MASI à inclure (par secteur,
    liquidité ou valeur de signal de consensus)~;
    \item \textbf{Capital}~: budget total et currency~;
    \item \textbf{Allocation par action}~: méthode d'allocation (HRP, Markowitz)~;
    \item \textbf{Entrées/Sorties et Dimensionnement}~: règles d'entrée, de sortie,
    sizing (\texttt{size\_pct}, Kelly), ATR stop~;
    \item \textbf{Risque}~: stop-loss global, drawdown maximum, cooldown.
\end{enumerate}

\subsection{Déclenchement du mode WFO}

Chaque paramètre peut être marqué \texttt{optimize=True}. Lorsqu'au moins un paramètre
est ainsi marqué, la page Backtest détecte automatiquement le mode WFO et affiche les
contrôles de configuration supplémentaires (dates WFO, horizon, période test).

\figuretodo{screenshot_strategy_page}{Capture d'écran de la page Stratégie~: 5 blocs dépliables, paramètre fenêtre SMA avec toggle "Optimiser via WFO" activé, presets short/medium/long.}

\section{Station 4 --- Backtest}
\label{sec:station_backtest}

\subsection{Deux modes de fonctionnement}

\begin{itemize}
    \item \textbf{Mode Standard}~: backtest sur une période unique avec paramètres
    fixés manuellement. Utile pour les évaluations rapides ou les stratégies à
    paramètres connus.
    \item \textbf{Mode WFO}~: déclenché automatiquement si des paramètres sont
    marqués \texttt{optimize=True}. Implémente le pipeline complet du
    Chapitre~\ref{ch:wfo}.
\end{itemize}

\subsection{Modèle de coût réaliste}

Le backtest intègre un modèle de coût de transaction calibré sur les conditions
de marché marocaines~:

\begin{itemize}
    \item \textbf{Frais de courtage}~: en points de base du montant de la transaction~;
    \item \textbf{Commission}~: frais fixes ou proportionnels~;
    \item \textbf{Slippage}~: estimation de l'impact de marché~;
    \item \textbf{TVA}~: taxe sur la valeur ajoutée applicable aux services
    financiers au Maroc.
\end{itemize}

\figuretodo{screenshot_backtest_page}{Capture d'écran de la page Backtest en mode WFO~: sélecteur de stratégie, contrôles WFO (dates, horizon, période test), barre de progression, résultats WFE par titre.}

\section{Persistance, Traçabilité et Bonnes Pratiques}
\label{sec:persistance}

\subsection{Module d'artefacts canonique}

Un module partagé \texttt{core/quant\_core/artifacts/} centralise toutes les
interactions avec le stockage objet~:

\begin{itemize}
    \item \texttt{errors.py}~: exceptions typées (\texttt{ArtifactNotFound},
    \texttt{StorageAuthError}, \texttt{StorageUnavailable})~;
    \item \texttt{models.py}~: dataclass \texttt{DatasetRef}~;
    \item \texttt{store.py}~: \texttt{head\_object}, \texttt{get\_bytes},
    \texttt{download\_to\_temp}~;
    \item \texttt{datasets.py}~: \texttt{get\_dataset\_ref}, \texttt{preflight\_dataset},
    \texttt{ensure\_dataset\_materialized}.
\end{itemize}

Cette centralisation garantit qu'API et Worker utilisent exactement le même code
pour accéder au stockage — pas de dérive entre implémentations.

\subsection{Tests automatisés}

\begin{itemize}
    \item 87+ tests unitaires et d'intégration (\texttt{pytest})~;
    \item Tests de régression sur le chemin d'optimisation rapide~;
    \item Tests de format et de cohérence des datasets.
\end{itemize}

\section*{Conclusion du chapitre}

Ce chapitre a présenté l'architecture technique qui héberge et orchestre les deux
moteurs de résultat. La séparation stricte en quatre stations, le stack full-stack
industrialisable et la traçabilité complète des runs constituent la fondation sur
laquelle repose le tableau de bord final, présenté dans le chapitre suivant.
